\section{Related Work}
\label{sec:related}

\noindent\textbf{Minecraft Agents} are intelligent agents designed to accomplish various tasks while playing the game Minecraft. Most of previous works adopt reinforcement learning for agent training. Due to the extremely sparse rewards and complex decision space of Minecraft tasks, early attempts have tried hierarchical RL~\cite{tessler2017deep,mao2022seihai,lin2022juewu,skrynnik2021hierarchical}, curriculum learning~\cite{kanitscheider2021multi}, and imitation learning~\cite{amiranashvili2020scaling} to empower more efficient RL training. To narrow the decision space, recent work~\cite{baker2022video} build a foundation model by performing imitation on YouTube videos. DreamerV3~\cite{hafner2023mastering} instead learns a world model that explores the environment efficiently. As the LLMs demonstrate their general planning ability, a series of research~\cite{huang2022inner,wang2023describe,yuan2023plan4mc,zhu2023ghost,wang2023voyager} leverage LLMs as high-level planners that decompose long-term complex tasks as basic skills and implement the skills with RL agents or handcrafted scripts.

Auto MC-Reward aims to design dense rewards for Minecraft tasks automatically using LLMs, which is orthogonal to previous works on Minecraft agents that mainly focus on RL learning algorithms or high-level planning.

\vspace{0.5em}\noindent\textbf{Efficient Learning in Sparse Reward Tasks} is a long-standing challenge in RL due to the lack of effective learning signals~\cite{ladosz2022exploration}. A common solution is to handcraft dense reward functions that provide intermediate reward signals based on human expertise, which requires time-consuming trial-and-error for each environment and task. Another line of previous research focus on proposing general-purpose dense auxiliary reward functions, such as curiosity-driven exploration~\cite{jaderberg2016reinforcement,pathak2017curiosity,burda2018large,li2020random}, self-imitation learning~\cite{oh2018self}, and goal-conditioned reinforcement learning~\cite{andrychowicz2017hindsight,levy2017learning,kulkarni2016hierarchical,vezhnevets2017feudal}. Despite the success on certain specific tasks, the applicability of these methods in the complex environment of Minecraft remains uncertain. 
Recent works~\cite{fan2022minedojo,ma2022vip,ma2023liv,du2023vision,du2023guiding} also propose to use pre-trained models to assign reward to intermediate states of completing tasks. However, these approaches produce reward values in a black-box manner, which cannot be interpreted and improved based on the experience of the agents, and the generalizability of these models on new tasks is not guaranteed.

In contrast, Auto MC-Reward automatically produces explainable reward functions according to the task descriptions. Moreover, the reward functions can be improved to be more precise based on the experience of the agent.


\vspace{0.5em}\noindent\textbf{Automated Reward Function Design} aims to find an optimal reward function that drives efficient reinforcement learning for interested tasks. Previous works~\cite{chiang2019learning,faust2019evolving} employ evolutionary algorithm for searching optimal reward functions for specific tasks. Most of these attempts have a highly constrained search space that only adjusts parameters of task-specific handcrafted reward templates. Recently, a series of research~\cite{du2023guiding,kwon2023reward,lin2022inferring,carta2022eager} employs LLMs for integrating human preference into open-domain tasks without clear completion criteria by directly prompting LLMs with environment trajectories and natural language task descriptions. The reward values are generated on-the-fly by LLMs, which is black-box and has heavy computational cost due to the nature of LLMs. In contrast, Auto MC-Reward employs LLMs to generate white-box code-form reward functions.

Concurrent works~\cite{yu2023language,xie2023text2reward,ma2023eureka} also propose to use LLMs as a coder to generate reward functions for robotics control tasks. Specifically, L2R~\cite{yu2023language} needs to prepare reward function templates and cannot cope with unexpected situations in open worlds. Text2Reward~\cite{xie2023text2reward} and EUREKA~\cite{ma2023eureka} require complete environment code or description and rely on human feedback, which are not available in open worlds. Different from these methods, Auto MC-Reward considers more complex Minecraft environments that has diverse scenarios and high uncertainty, requiring more precise and thorough reward designing.
